\subsection{A one-edge circle against an arbitrary circle graph}

\begin{theorem}[The one-edge-circle chart obstruction]
\label{C:one-edge-circle-chart-zero}
Use the corner monoids and three-object support of
Theorem~\ref{C:affine-chart-family-zero}.  Suppose one adjacent
act is the two-state act whose two retractions both send its
transient state to its fixed state.  Suppose the natural graph
of the other adjacent act is connected and has first Betti
number one.  The long-edge biset and composition may be arbitrary.
Then every rational product of degree-two classes on $BC$ is
zero.  The same assertion holds with the two adjacent acts
interchanged.
\end{theorem}
\begin{proof}
Let $A$ be the one-edge circle and $B$ the other action graph.
Its unique essential circuit is a graph deformation retract.
Restricting to that circuit in the triangle-of-spaces model
is an objectwise homotopy equivalence: the face becomes
$A\times B_0\to A\times B$, where $B_0$ is that circuit,
and the long edge stays unchanged.  We may therefore assume
that $B$ is a circle with any finite number of oriented edges.
Write $\varepsilon_e\in\{1,-1\}$ for their coefficients in
an oriented fundamental cycle.

Use the natural cubical model $K=K(T)$ from
Lemma~\ref{C:two-retraction-cubical-model}, and write
\[
 h:A\times B\longrightarrow K
\]
for composition.  For an edge $e$ of $B$, let $t_e$ be the
long-edge state assigned to the product of $e$ with the unique
edge of $A$.  The two endpoints of that $A$-edge are the same
fixed state.  Consequently
\[
 P_0t_e=P_1t_e.
\]
As before, the cellular image of the torus fundamental class is
\[
 z=\sum_e\varepsilon_e[t_e]\in C_2(K;\mathbb Z),
\]
where $[t_e]$ is zero for a degenerate cell.  If $z\ne0$,
then $h_*:H_2(A\times B;\mathbb Q)\to H_2(K;\mathbb Q)$
is injective, since $K$ has no three-cells.  The triangle
spectral sequence gives $H_4(BC;\mathbb Q)=0$, proving the
claim in this case.

Suppose $z=0$.  If some $t_e$ belongs to $U$, the entire
$A$-coordinate collapses over that edge.  In particular the
$A$-loop at an endpoint of $e$ maps to a point.  Transporting
the basepoint along $B$ shows that its image is null-homotopic
at every choice of $B$-basepoint.  Thus
\begin{equation}\label{C:loop-chart-one-coordinate}
 h_*H_1(A;\mathbb Z)=0.
\end{equation}

Otherwise no $t_e$ on the essential circuit belongs to $U$.
At the fixed vertex of $A$, the image of the $B$-edge $e$
is the one-cell indexed by $P_0t_e$, or a vertex when that
state lies in $U\cap V$.  If $t_e\in V$, commutation makes
$P_0t_e\in U\cap V$, so its contribution is zero.  Every
remaining $t_e$ is a nondegenerate two-cell.  Apply to the
literally zero chain $z$ the linear map
\[
 C_2(K;\mathbb Z)\longrightarrow C_1(K;\mathbb Z),
 \qquad [t]\longmapsto
 \begin{cases}
 [P_0t],&P_0t\in U\setminus V,\\
 0,&P_0t\in U\cap V.
 \end{cases}
\]
Its value is exactly the cellular image of the $B$-cycle.
Hence in this second case
\begin{equation}\label{C:loop-chart-other-coordinate}
 h_*H_1(B;\mathbb Z)=0.
\end{equation}
We do not claim that the full map $h$ factors through a
projection: only one of these two homology statements is
needed.

Here is the cohomological consequence, which also explains
why homology in the preceding step suffices.  The triangle
spectral sequence identifies
\[
 H^2(BC;\mathbb Q)\cong
 \ker\bigl(H^1(A;\mathbb Q)\oplus H^1(B;\mathbb Q)
             \oplus H^1(K;\mathbb Q)
           \longrightarrow H^1(A\times B;\mathbb Q)\bigr).
\]
Projection of this kernel onto $H^1(K;\mathbb Q)$ is an
isomorphism, because the first two edge coordinates map
isomorphically onto the two summands of the target.  The
edge-to-face map in degree zero is onto, so there is no
additional degree-two term.  This remains valid when $K$
is disconnected.

For any integral class $u\in H^1(K;\mathbb Z)$, choose a
representing map $f_u:K\to S^1$.  Replacing the long edge
by $S^1$ and its face map by $f_u h$ gives a triangle of
spaces $X_u$ and an actual map
\[
 BC\longrightarrow X_u.
\]
The generator of $H^2(X_u;\mathbb Q)$ pulls back to the
class corresponding to $u$ in the preceding kernel
identification.  By \eqref{C:loop-chart-one-coordinate} or
\eqref{C:loop-chart-other-coordinate}, the map
$f_u h:A\times B\to S^1$ has degree zero on one coordinate.
Since $S^1$ is a $K(\mathbb Z,1)$, it is homotopic to a map
through the other projection.  The suspension construction
of Theorem~\ref{C:affine-chart-family-zero} then supplies a
map $X_u\to\Sigma S^1$ inducing an isomorphism on degree-two
cohomology.
The generator on $X_u$ therefore has zero square, and so
does its pullback on $BC$.

The finite CW complex $K$ has an integral first-cohomology
lattice spanning $H^1(K;\mathbb Q)$.  Multiplying any rational
class by an integer reduces it to the preceding integral
case.  Thus every rational degree-two square on $BC$ is
zero.  Polarization gives the vanishing of every product
of two such classes.
\end{proof}

\begin{corollary}[A one-edge circle against any connected action graph]
\label{C:one-edge-arbitrary-graph-cups}
In Theorem~\ref{C:one-edge-circle-chart-zero}, the other adjacent
action graph may be any connected finite graph, with no restriction
on its first Betti number.  All rational degree-two products still
vanish.
\end{corollary}
\begin{proof}
Write $X$ for the same triangle-of-spaces model, with adjacent
edges $A=S^1$ and $B$, long edge $K=K(T)$, and composition
$h:A\times B\to K$.  All edge spaces have dimension at most two,
the vertex spaces are points, and the face has dimension two.
The integral triangle spectral sequence therefore identifies
\[
 H_4(X;\mathbb Z)=
 \ker\bigl(H_2(A\times B;\mathbb Z)
                 \xrightarrow{h_*}H_2(K;\mathbb Z)\bigr).
\]
There are no other terms in total degree four and no higher
differential involving this kernel.  Since
$H_2(A\times B;\mathbb Z)=H_1(A;\mathbb Z)\otimes H_1(B;\mathbb Z)$,
each element of this kernel has the form $[A]\otimes\gamma$.

Choose a based closed edge word $\lambda:S^1\to B$ representing
$\gamma$.  This is always possible by choosing a maximal tree
in the connected graph and writing a word in its remaining edges.
Pulling back the triangle along $\lambda$ gives a triangle
$X_\gamma$ and an actual map $X_\gamma\to X$.  Its face map is
$h_\gamma=h(1\times\lambda):A\times S^1\to K$.
The cellular image of its fundamental two-cycle is literally
zero: it represents $h_*([A]\otimes\gamma)=0$, and $K$ has
no three-cells.  Consequently $H_4(X_\gamma;\mathbb Z)$ has
a generator mapping to the specified element of $H_4(X;\mathbb Z)$.

The zero-chain part of the preceding proof applies verbatim
to this closed edge word.  More explicitly, each traversal
copies the corresponding product square, with its orientation
reversed when the edge is traversed backwards.  If any used
top state lies in $U$, the $A$-loop has constant image at a
point of the other circle and hence zero image in first homology.
Otherwise every used degenerate state lies in $V$, and applying
the displayed $P_0$-face map to the literally zero signed
two-chain shows that the other circle has zero image in first
homology.  Repeated edges merely repeat their signed summands;
no embedded-circuit assumption is used in either argument.
For every map $K\to S^1$, the resulting face map therefore has
one coordinate degree zero.  The same suspension comparison
proves that every degree-two square on $X_\gamma$ is zero.
It is not necessary for $X_\gamma$ itself to be the classifying
space of a category in the stated family.

By naturality, the square of any rational degree-two class on
$X$ evaluates to zero on every integral element of $H_4(X)$.
These integral kernel elements span $H_4(X;\mathbb Q)$, and
rational cohomology is detected by evaluation on rational
homology.  Every rational square is therefore zero; polarization
proves the assertion for all mixed products.
\end{proof}

\begin{remark}[The complete longer-circle controls]
\label{C:longer-circle-chart-controls}
The exact program \path{code/d_circle_charts.py}, with the
typed filter \path{code/d_circle_charts_filter.cpp}, also
checks the following two bounded families.  Circle orientations
are retained explicitly: a two-edge circuit is either two
equally directed parallel edges or a directed two-cycle;
a three-edge circuit is either cyclically directed or
transitively directed.
\[
\begin{array}{c|c|c|r}
\text{adjacent state counts}&|T|&\text{orientation pairs}
 &\text{checked candidates}\\\hline
(4,6)&1,\ldots,4&4&5844\\
(2,6)&1,\ldots,6&2&44304.
\end{array}
\]
Every commuting long-biset retraction pair and every equivariant
composition is included.  The only orbit reductions are actual
state automorphisms of the two adjacent acts and the long biset.
The completed independent bar checks give explicit zero-square
primitives for all $50148$ candidates.  The output prefixes
\path{code/outputs/d_circle_charts_} record the chosen graph
names and long-biset size; \texttt{parallel2}, \texttt{cycle2},
\texttt{cycle3}, and \texttt{acyclic3} specify these four
orientations.  The additional suffix \texttt{minimal} denotes
a separately restricted fixed four-state target, not the full
biset search.  The theorem above makes the second row an
all-size structural exclusion.  The first row is stated only
in its completed bounded scope.
\end{remark}
